\documentclass[10pt]{beamer}
\usetheme[
  ]{Feather}

\usepackage[utf8]{inputenc}
\usepackage[english,brazil]{babel}
\usepackage[T1]{fontenc}
\usepackage{helvet}
\usepackage{hyperref}
\usepackage{tabularx}
\usepackage{lmodern}
\usepackage{float}
\usepackage{threeparttable}
\usepackage{array}
\usepackage{booktabs}
\usepackage{caption}
\usepackage{tikz}
\usetikzlibrary{trees, shapes.geometric, positioning, arrows.meta, calc}


\title[]% [] is optional
{%
  \textbf{Geração Automatizada de Testes Baseada em Model Checking:
  \\ Um Estudo Aplicado ao e-Título}
}

\subtitle[Geração Automatizada de Testes Baseada em Model Checking]
{
  \textbf{Versão para Qualificação}
}

\author[Thiago Fernandes]
{Thiago Viana Fernandes \\
  {\ttfamily thiagu.vf@gmail.com}
}

\institute[]
{
  Programa de Pós-Graduação em Computação Aplicada\\
  Universidade de Brasília
}

\date{\today}

\begin{document}

% Title page
{%
\begin{frame}[plain,noframenumbering]
  \titlepage
\end{frame}}

\section{Introdução}
\begin{frame}{Motivação}{Contexto}
  \begin{itemize}
    \item APIs REST são base de sistemas modernos e críticos.
    \item Erros de integração frequentemente não são detectados por testes tradicionais.
    \item Testes manuais e exploratórios possuem baixa sistematização.
    \item Há lacuna entre modelos formais e práticas de teste de software.
  \end{itemize}
\end{frame}

\begin{frame}{Problema}{Desafio}
  \begin{itemize}
    \item Como gerar testes de integração de forma sistemática e baseada em especificação?
    \item Como garantir coerência entre regras de negócio, serviços REST e testes?
    \item Como explorar contra-exemplos de Model Checking além da verificação?
  \end{itemize}
\end{frame}

\begin{frame}{Ideia Central}{Visão Geral}
  \begin{itemize}
    \item Modelar comportamento do sistema.
    \item Verificar propriedades com Model Checking.
    \item Usar contra-exemplos como fonte de contratos e testes.
  \end{itemize}
\end{frame}



\begin{frame}{Fundamentação}{Model Checking}
  \begin{itemize}
    \item Técnica formal para verificação automática de propriedades.
    \item Baseada em exploração exaustiva de estados.
    \item Propriedades expressas em LTL.
    \item Ferramenta utilizada: SPIN.
  \end{itemize}
\end{frame}

\begin{frame}{Fundamentação}{Design by Contract}
  \begin{itemize}
    \item Pré-condições.
    \begin{figure}[hb]
	    \resizebox{0.3\textwidth}{!}{
        \begin{tikzpicture}
  [state/.style={circle,draw=black,minimum size=8mm},
   guard/.style={align=center,font=\itshape}]

  % Estados
  \node(s0) at (0,0) [state] {$s_0$};
  \node(s1) at (3,0) [state] {$s_1$};

  % Transição com pré-condição
  \draw [-{Stealth[round]}] (s0) to node[above,guard] {$[x > 0]$} (s1);

\end{tikzpicture}
      }
    \end{figure}
    \item Pós-condições.
    \begin{figure}[hb]
	    \resizebox{0.3\textwidth}{!}{
        \begin{tikzpicture}
  [state/.style={circle,draw=black,minimum size=8mm},
   post/.style={align=center,font=\itshape}]

  % Estados
  \node(s0) at (0,0) [state] {$s_0$};
  \node(s1) at (3,0) [state] {$s_1$};

  % Transição
  \draw [-{Stealth[round]}] (s0) to (s1);

  % Pós-condição
  \node at (s1.south) [post,below] {$\{y = 1\}$};

\end{tikzpicture}
      }
    \end{figure}
    \item Invariantes.
    \begin{figure}[hb]
	    \resizebox{0.3\textwidth}{!}{
        \begin{tikzpicture}
  [state/.style={circle,draw=black,minimum size=8mm},
   inv/.style={align=center,font=\itshape}]

  % Estados
  \node(s0) at (0,0) [state] {$s_0$};
  \node(s1) at (3,0) [state] {$s_1$};

  % Invariante nos estados
  \node at (s0.south) [inv,below] {$\{z \geq 0\}$};
  \node at (s1.south) [inv,below] {$\{z \geq 0\}$};

  % Transição
  \draw [-{Stealth[round]}] (s0) to (s1);

\end{tikzpicture}
      }
    \end{figure}
    \item Contratos como oráculos de teste.
  \end{itemize}

\end{frame}

\begin{frame}{Fundamentação}{Testes de Integração}
  \begin{itemize}
    \item Validação de contratos da API.
    \item Sequências de chamadas.
    \item Estados intermediários.
    \item OpenAPI como suporte à especificação.
  \end{itemize}
\end{frame}

\begin{frame}{Método Proposto}{Visão Geral}
  \begin{itemize}
    \item Modelagem comportamental em Promela.
    \item Definição de propriedades LTL.
    \item Execução do Model Checker.
    \item Extração de contra-exemplos.
    \item Geração de contratos e testes.
  \end{itemize}

\end{frame}

\begin{frame}{Método Proposto}{Visão Geral}
  \begin{center}
    \begin{figure}[hb]
	    \resizebox{0.9\textwidth}{!}{
        \begin{tikzpicture}
    [artefato/.style={rectangle,draw=black, thick, inner sep=10pt, minimum size=15mm, rounded corners, align=center},
     atividade/.style={circle, draw=black, thick, inner sep=10pt, minimum size=15mm, align=center}]

    % node (name) at (,) [type]{label};
    \node[artefato] (def_alt_niv) at(0,6)  {Definição em \\alto nível};
    \node[artefato] (cod_api)     at(0,3)  {Código Promela \\API};
    \node[artefato] (mod_form)    at(0,0)  {Modelo Formal \\Compilado};
    \node[atividade](valid)       at(5,0)  {Verificação \\propriedades \\LTL};
    \node[artefato] (claus_dbc)   at(11,0) {Cláusulas \\DBC};

    \node[above] at (valid.south) {$\phi = \cdots$};

    \draw [-{Stealth}, thick] (def_alt_niv) -- node[auto, swap] {\textit{interpretação}} (cod_api);
    \draw [-{Stealth[green!200]}, thick, color=green] (cod_api) -- node[auto, swap]{\textit{modelagem}} (mod_form);
    \draw [-{Stealth}, thick] (mod_form) -- node[auto]{\textit{verificação}} (valid);
    \draw [-{Stealth[orange]}, thick, color=orange] (valid) -- node[auto]{\textit{extr. cláusulas}} (claus_dbc);
    \draw [-{Stealth}, thick] (valid) to [bend left=45] node [auto]{\textit{ajustes modelo}} (mod_form);
    \draw [-{Stealth}, thick] (valid) to [bend right=20] node[auto, swap]{\textit{ajustes código}} (cod_api);
    \draw [-{Stealth[red]}, thick, color=red] (claus_dbc) to [bend right=30] node[auto, swap]{\textit{especificações DBC}} (cod_api);

\end{tikzpicture}

      }
    \end{figure}
  \end{center}
\end{frame}

\begin{frame}{Método Proposto}{Modelagem}
  \begin{itemize}
    \item Processos representando serviços.
    \item Estados explícitos.
    \item Comunicação por variáveis globais.
  \end{itemize}
\end{frame}

\begin{frame}{Método Proposto}{Propriedades}
  \begin{itemize}
    \item Especificam comportamentos esperados.
    \item Exemplo: após carrossel válido, senha deve ser definida.
  \end{itemize}

\begin{center}
  \texttt{[] (carrossel\_ok -> <> senha\_definida)}
\end{center}

\end{frame}

\begin{frame}{Método Proposto}{Contra-exemplo}
  \begin{itemize}
    \item Caminho que viola propriedade.
    \item Sequência concreta de estados.
    \item Evidência de falha lógica.
  \end{itemize}
\end{frame}

\begin{frame}{Método Proposto}{Derivação de Contratos}
  \begin{itemize}
    \item Estados iniciais → pré-condições.
    \item Estado final inválido → pós-condição.
    \item Regras como cláusulas formais.
  \end{itemize}
\end{frame}

\begin{frame}{Método Proposto}{Geração de Testes}
  \begin{itemize}
    \item Pré-condições → dados de entrada.
    \item Sequência → passos do teste.
    \item Pós-condições → asserções.
  \end{itemize}
\end{frame}

\begin{frame}{Estudo de Caso}{Aplicativo e-Título}
  \begin{itemize}
    \item Fluxo de onboarding.
    \resizebox{0.8\textwidth}{!}{
	    \begin{tikzpicture}
    [initial/.style={circle,draw=black,fill=black,thick}, 
     state/.style={circle,draw=black,fill=white,thick, minimum size=11mm},
     decision/.style={diamond,draw=black,fill=white,thick, minimum size=6mm},
     final/.style={circle,draw=black, fill=black, thick, double}]
     \node[initial] (inicial) {};
     \node[state] (dados_cadastrais) [right=of inicial] {$s_1$};
     \node[state] (carrossel) [right=of dados_cadastrais] {$s_2$};
     \node[state] (user_senha) [right=of carrossel] {$s_3$};
     \node[decision, align=center] (quest_biometria) [right=of user_senha] {\small Biometria? \\ $s_5$};
     \node[state] (biometria) [right=of quest_biometria] {$s_6$};
     \node[decision] (aprovado) [right=of biometria] {$s_7$};
     \node[state] (etituloEmitido) [right=of aprovado] {$s_8$};
     \node[state] (bloqueado) [below=of carrossel] {$s_4$};
     \node[state] (etituloReduzido) [below=of etituloEmitido] {$s_9$};
     \node[final] (final) [right=of etituloEmitido] {};

     % nodos de texto
     \node [below, align=center] at (inicial.south) {$s_0$};
     \node [above, align=center] at (dados_cadastrais.north) {Dados\\ cadastrais};
     \node [above, align=center] at (carrossel.north) {Carrossel};
     \node [below, align=center] at (bloqueado.south) {Bloqueado};
     \node [above, align=center] at (user_senha.north) {Usuário\\e senha};
     \node [above, align=center] at (quest_biometria.north) {};
     \node [above, align=center] at (biometria.north) {Realizar\\biometria};
     \node [above, align=center] at (etituloEmitido.north) {e-Titulo\\emitido};
     \node [below, align=center] at (etituloReduzido.south) {e-Título\\reduzido};
     \node [below, align=center] at (final.south) {$s_{10}$};
     
     % arestas
     \draw[-{Stealth[round]}] (inicial) to [thick] (dados_cadastrais); 
     \draw[-{Stealth[round]}] (dados_cadastrais) to [thick] (carrossel);
     \draw[-{Stealth[round]}] (carrossel) to [thick] (user_senha);
     \draw[-{Stealth[round]}] (carrossel) to [thick] (bloqueado);
     \draw[-{Stealth[round]}] (user_senha) to [thick] (quest_biometria);
     \draw[-{Stealth[round]}] (quest_biometria) to [thick] node[auto,swap]{\textit{sim}}(biometria);
     \draw[-{Stealth[round]}] (biometria) to [thick] (aprovado);
     \draw[-{Stealth[round]}] (aprovado) to [thick] node[auto]{\tiny \textit{aprovado}} (etituloEmitido);
     \draw[-{Stealth[round]}] (etituloEmitido) to [thick] (final);
     \draw[-{Stealth[round]}] (etituloReduzido) to [thick] (final);
     \draw[-{Stealth[round]}] (bloqueado) to [thick, bend left] (dados_cadastrais);
     \draw[-{Stealth[round]}] (quest_biometria) to [thick, out=270, in=180] node[auto,swap] {\textit{não}} (etituloReduzido);
     \draw[-{Stealth[round]}] (aprovado) to [thick, out=270, in=135] node[auto,swap]{\small\textit{reprovado}} (etituloReduzido);
\end{tikzpicture}

    }
    \item Carrossel, senha e biometria.
    \item Sistema real e crítico.
  \end{itemize}
\end{frame}

\begin{frame}{Resultados}{Síntese}
  \begin{itemize}
    \item Propriedades violadas identificadas.
    \item Contratos derivados automaticamente.
    \item Casos de teste obtidos.
    \item Detecção de cenários não cobertos.
  \end{itemize}
\end{frame}

\begin{frame}{Conclusão}{Síntese}
  \begin{itemize}
    \item Proposta de integração entre Model Checking, Design by Contract e testes de APIs REST.
    \item Demonstração de que contra-exemplos podem ser sistematicamente convertidos em cláusulas contratuais e casos de teste.
    \item Evidência de que a abordagem auxilia na identificação de falhas lógicas e de uso indevido de serviços.
    \item Validação em um fluxo real do aplicativo e-Título.
  \end{itemize}
\end{frame}

{
\begin{frame}[plain,noframenumbering]
  \finalpage{Obrigado!}
\end{frame}}

\end{document}
